\subsection{Sub-Phase 3.1: State-Dependent Tangent Stiffness Assembly}
\label{subsec:tangent_stiff}

To model structural components undergoing substantial structural deflection safely, a geometrically nonlinear formulation was integrated into the continuous 2D element loop. Under large displacement conditions, the linear strain assumptions become invalid, requiring the introduction of the Green-Lagrange strain tensor $\mat{E}$:
\begin{equation}
    \mat{E} = \frac{1}{2} \left( \nabla \vect{u} + \nabla \vect{u}^T + \nabla \vect{u}^T \nabla \vect{u} \right)
\end{equation}
where $\nabla \vect{u}$ represents the spatial gradient of the continuous displacement field evaluated via the bivariate NURBS shape functions. 

Consequently, the internal force vector $\vect{F}_{\text{int}}(\vect{u})$ and the state-dependent tangent stiffness matrix $\mat{K}_T(\vect{u})$ must be continuously re-evaluated. The internal force vector is computed over the reference configuration $\Omega_0$ using the nonlinear strain-displacement matrix $\mat{B}_L(\vect{u})$ and the second Piola-Kirchhoff stress vector $\vect{S} = [S_{xx}, S_{yy}, S_{xy}]^T$:
\begin{equation}
        \vect{F}_{\text{int}}(\vect{u}) = \int_{\Omega_0} \mat{B}_L^T(\vect{u}) \vect{S} \, d\Omega_0
\end{equation}
where the nonlinear strain-displacement matrix is split into a displacement-independent part $\mat{B}_0$ and a displacement-dependent nonlinear part $\mat{B}_{\text{NL}}(\vect{u})$:
\begin{equation}
        \mat{B}_L(\vect{u}) = \mat{B}_0 + \mat{B}_{\text{NL}}(\vect{u})
\end{equation}
The continuous tangent stiffness operator $\mat{K}_T(\vect{u})$ is obtained by taking the variation of the internal force vector, which decomposes into elastic (material) and geometric (initial stress) components:
\begin{equation}
    \mat{K}_T(\vect{u}) = \mat{K}_e(\vect{u}) + \mat{K}_{\sigma}(\vect{u}) = \left( \mat{K}_0 + \mat{K}_L(\vect{u}) \right) + \mat{K}_{\sigma}(\vect{u})
\end{equation}
Here, $\mat{K}_0$ represents the standard baseline material stiffness matrix, whereas the large displacement kinematic stiffness matrix $\mat{K}_L(\vect{u})$ accounts for the kinematic coupling introduced by finite deflections, formulated as:
\begin{equation}
        \mat{K}_L(\vect{u}) = \int_{\Omega_0} \left( \mat{B}_0^T \mat{D} \mat{B}_{\text{NL}}(\vect{u}) + \mat{B}_{\text{NL}}^T(\vect{u}) \mat{D} \mat{B}_0 + \mat{B}_{\text{NL}}^T(\vect{u}) \mat{D} \mat{B}_{\text{NL}}(\vect{u}) \right) \, d\Omega_0
\end{equation}
The initial stress geometric stiffness matrix $\mat{K}_{\sigma}(\vect{u})$ accounts for structural stiffness changes due to internal stresses and is defined by:
\begin{equation}
        \mat{K}_{\sigma}(\vect{u}) = \int_{\Omega_0} \mat{G}^T \mat{M}(\vect{S}) \mat{G} \, d\Omega_0
\end{equation}
For each control point node $a$, the displacement gradient interpolation block $\mat{G}_a \in \mathbb{R}^{4 \times 2}$ is given by:
\begin{equation}
        \mat{G}_a = \begin{bmatrix} \frac{\partial R_a}{\partial x} & 0 \\[0.3em] \frac{\partial R_a}{\partial y} & 0 \\[0.3em] 0 & \frac{\partial R_a}{\partial x} \\[0.3em] 0 & \frac{\partial R_a}{\partial y} \end{bmatrix}
\end{equation}
and the stress tensor block matrix $\mat{M}(\vect{S}) \in \mathbb{R}^{4 \times 4}$ maps the components of the second Piola-Kirchhoff stress:
\begin{equation}
        \mat{M}(\vect{S}) = \begin{bmatrix} \mat{S}_{\text{mat}} & \mat{0} \\ \mat{0} & \mat{S}_{\text{mat}} \end{bmatrix}, \quad \text{with} \quad \mat{S}_{\text{mat}} = \begin{bmatrix} S_{xx} & S_{xy} \\ S_{xy} & S_{yy} \end{bmatrix}
\end{equation}

\begin{figure}[H]
    \centering
    % [IMAGE PLACEHOLDER: State-dependent tangent operator layout or Jacobian structure]
    \vspace{4cm}
    \caption{Jacobian sparsity pattern and matrix profile adjustments during high-order state configurations within the continuous 2D plane stress elements.}
    \label{fig:jacobian_pattern}
\end{figure}

The evaluation core maps the state-dependent strain-displacement matrix $\mat{B}_N(\vect{u})$ directly to the Gauss-Legendre integration points, tracking the incremental energy change across each element span.

\begin{devnote}
The mathematical assembly sequence mapped inside \texttt{phase-3-core/src/iga-nonlinear.js} isolates the calculation of the geometric matrix $\mat{K}_{\sigma}$ from the standard material operators. This isolation permits multi-threaded caching transformations if the structural domain needs to be accelerated via Web Workers or similar pipeline parallelisms during rapid runtime execution loops.
\end{devnote}